\subsection{Current Consumption Analysis}
\cref{fig:cc_full} shows the measured current consumption of the full system 
across a temperature range of \SI{-20}{\celsius} to \SI{+50}{\celsius} at constant 
humidity. The two operating modes exhibit opposing temperature dependencies, each 
governed by distinct physical mechanisms.

\begin{figure}[H]
\centering
\begin{subfigure}{0.48\textwidth}
    \centering
    \includestandalone[mode=buildmissing, width=\textwidth]{Standalone/Plots/FullSystem_neg_deg_cc_meas}
    \caption{EM0 (run mode)}
    \label{fig:measurement}
\end{subfigure}
\hfill
\begin{subfigure}{0.48\textwidth}
    \centering
    \includestandalone[mode=buildmissing, width=\textwidth]{Standalone/Plots/FullSystem_neg_deg_cc_sleep}
    \caption{EM4 (sleep mode)}
    \label{fig:sleep}
\end{subfigure}
\caption{Full system current consumption of EM0 (run mode) and EM4 (sleep mode).}
\label{fig:cc_full}
\end{figure}

In EM0 (run mode), the average current consumption increases at lower temperatures,  from \SI{5.612}{\milli\ampere} at \SI{+50}{\celsius} to \SI{6.681}{\milli\ampere} at \SI{-20}{\celsius}, with a total average at \SI{5.886}{\milli\ampere} for the whole time period.This trend is opposite to the intrinsic MCU behavior characterized in the EFR32MG22E datasheet \cite{EFR32MG22E}, which shows EM0 supply current increasing with temperature when powered from a stable regulated supply. The reason for this reversal can be found in the temperature dependence of the EVE ER14250 Li-SOCl\textsubscript{2} battery. The output of the battery was measured as shown in \cref{fig:FS_bat}. Here the battery terminal voltage drops significantly at low temperatures with active-pulse voltages approaching \SI{2.6}{\volt} at \SI{-20}{\celsius}. 

\begin{figure}[H]
\centering
\includestandalone[mode=buildmissing, width=0.75\textwidth]{Standalone/Plots/FullSystem_neg_deg_bat}
    \caption{Battery voltage in varying enviromental conditions.}
    \label{fig:FS_bat}
\end{figure}

Since the MCU digital core, RF subsystem, and I/O supply rails (DVDD, RFVDD, PAVDD, AVDD, and IOVDD) are all powered by the integrated DC--DC buck converter operating at \(V_{\mathrm{DCDC}}=\SI{1.8}{\volt}\), the input current can be estimated as a function of the supply voltage under the assumption of a constant load. To investigate whether the observed decrease in input current with increasing supply voltage can be explained by the DC--DC converter, the run-mode load is modeled as a constant power demand \(P_{\mathrm{out}}\) supplied by the converter. The converter efficiency is assumed to be \(\eta \approx 0.91\), based on the efficiency characteristic provided in the datasheet \cite{EFR32MG22E}. The efficiency of a buck converter is defined as the ratio of output power to input power \cite[p.~93]{erickson2020fundamentals}:

\begin{equation}
    \eta = \frac{P_{\mathrm{out}}}{P_{\mathrm{in}}}.
\end{equation}

Using that the converter input power is given by \(P_{\mathrm{in}} = V_{\mathrm{in}} I_{\mathrm{in}}\) and the \SI{50}{\celsius} operating point (\(V_{\mathrm{in},50^\circ\mathrm{C}} \approx \SI{3.4}{\volt}\), \(I_{\mathrm{in},50^\circ\mathrm{C}} = \SI{5.612}{\milli\ampere}\)) as a reference, the delivered output power is

\begin{equation}
    P_{\mathrm{out}}
    = \eta V_{\mathrm{in},50^\circ\mathrm{C}} I_{\mathrm{in},50^\circ\mathrm{C}}
    = 0.91 \cdot \SI{3.4}{\volt} \cdot \SI{5.612}{\milli\ampere}
    = \SI{17.36}{\milli\watt}.
\end{equation}

Assuming that \(P_{\mathrm{out}}\) remains approximately constant with temperature, the reduced battery terminal voltage at \SI{-20}{\celsius} (\(V_{\mathrm{in},-20^\circ\mathrm{C}} \approx \SI{2.7}{\volt}\)) yields a predicted input current of

\begin{equation}
    I_{\mathrm{in},-20^\circ\mathrm{C}}
    = \frac{P_{\mathrm{out}}}{\eta V_{\mathrm{in},-20^\circ\mathrm{C}}}
    = \frac{\SI{17.36}{\milli\watt}}
           {0.91 \cdot \SI{2.7}{\volt}}
    \approx \SI{7.07}{\milli\ampere}.
\end{equation}

The predicted current of \SI{7.07}{\milli\ampere} is in good agreement with the measured value of \SI{6.681}{\milli\ampere}, corresponding to a deviation of approximately \(5.8\%\). This small discrepancy is expected, as the estimation assumes constant converter efficiency and a temperature-independent load power, while neglecting secondary effects such as temperature-dependent converter losses, variations in RF and digital subsystem power consumption, and battery impedance effects. Nevertheless, the result supports the hypothesis that the increase in measured input current at low temperatures is primarily caused by the reduction in battery terminal voltage, requiring the DC--DC converter to draw more current to maintain approximately constant output power.

In EM4 (sleep mode), the average current consumption rises with temperature, from \SI{171.9}{\nano\ampere} at \SI{-20}{\celsius} to \SI{767.7}{\nano\ampere} at \SI{+50}{\celsius}, with a total average of \SI{451}{\nano\ampere} for the whole time period. This trend matches the EFR32MG22E datasheet and confirms that sleep current dominates the long-term energy budget at low sample rates and elevated temperatures, as reflected in \autoref{tab:charge_10y}. It also proves that the microcontroller is operating at the lower power consumption possible, since that datasheet specified the absolute max at \SI{170}{\nano\ampere}. 

% The measured average current consumption in both operating modes at \SI{-20}{\celsius} and \SI{+50}{\celsius} is presented in \autoref{tab:cc_averages}. The results are consistent with the datasheet specifications, with the lowest measured EM4 current consumption being \SI{0.17}{\micro\ampere} and the EM0 current consumption ranging from \SI{5.612}{\micro\ampere} to \SI{6.674}{\micro\ampere} over the tested temperature range.

% \begin{table}[H]
% \centering
% \caption{Measured current consumption and charge per wake-up}
% \label{tab:cc_averages}
% \begin{tabular}{@{}lcr@{}}
% \toprule
% Mode &
% Temp. &
% $I_{\mathrm{avg}}$  \\
% &
% (\si{\celsius}) &
% (\si{\milli\ampere}) \\
% \midrule
% EM0 & -20 & \SI{6.674}{\milli\ampere} \\
% EM0 & +50 & \SI{5.612}{\milli\ampere} \\
% EM4 & -20 & \SI{171.9}{\nano\ampere} \\
% EM4 & +50 & \SI{767.7}{\nano\ampere} \\
% \bottomrule
% \end{tabular}
% \end{table}

A current consumption analysis is carried out to determine the nominal battery capacity required for a given number of reports over the device lifetime. Taking a \num{10}-year lifetime, the total operating time is
\begin{equation}
    H_{10\text{y}} = 24~\text{h/day} \times 365~\text{days/year}
        \times 10~\text{years} = \SI{87600}{\hour},
\end{equation}

The charge associated with a given current consumption is obtained from
\begin{equation}
    Q = I \cdot t,
\end{equation}
where $I$ is the current and $t$ the duration over which it is drawn.

In \cref{tab:charge_10y}, the average charge per cycle is used to determine the total nominal charge required to support different numbers of readings per day over a 10-year period. For the selected EVE ER14250 \SI{1200}{\milli\ampere\hour} battery, assuming a 10-year lifetime and an estimated capacity degradation of 1\% per year (resulting in a total loss of approximately 12\%, or \SI{144}{\milli\ampere\hour}), the maximum number of cycles per day is calculated to be approximately 55 at the highest current consumption at \SI{50}{\celsius}.

\begin{table}[H]
\centering
\caption{Estimated charge consumption over 10 years}
\label{tab:charge_10y}
\begin{tabular}{@{}crrrr@{}}
\toprule
Readings/day &
Temp. &
$Q_{\mathrm{EM0}}$ &
$Q_{\mathrm{EM4}}$ &
$Q_{\mathrm{total}}$ \\
Over 10 years  &
(\si{\celsius}) &
(\si{\milli\ampere\hour}) &
(\si{\milli\ampere\hour}) &
(\si{\milli\ampere\hour}) \\
\midrule
1 & -20 & 15.24 & 15.05 & 30.29 \\
1 & +50 & 18.12 & 67.23 & 85.35 \\
\hline
2 & -20 & 30.48 & 15.05 & 45.52 \\
2 & +50 & 36.24 & 67.21 & 103.45 \\
\hline
5 & -20 & 76.19 & 15.04 & 91.22 \\
5 & +50 & 90.61 & 67.15 & 157.75 \\
\hline
10 & -20 & 152.38 & 15.01 & 167.39 \\
10 & +50 & 181.21 & 67.04 & 248.25 \\
\hline
55 & -20 & 838.07 & 14.80 & 852.87 \\
55 & +50 & 996.67 & 66.10 & 1062.77 \\
\bottomrule
\end{tabular}
\end{table}

% In summary, based on the measured current consumption of the complete system implementation, it was demonstrated that the battery input voltage has the dominant influence on EM0 (run mode) current consumption. The experimental results follow the same overall trend as reported in the \cite{EFR32MG22E} datasheet, indicating good agreement between the measured behavior and the expected device characteristics.